\documentclass[10pt,a4paper]{amsart}
\usepackage[margin=19mm]{geometry}
\usepackage{amsmath,amssymb,mathtools}
\usepackage[scr=rsfs,cal=euler]{mathalfa}
\usepackage{xfrac}
\usepackage[unicode,hidelinks,bookmarks=false]{hyperref}
\newcommand{\ie}{i.e.\ }
\newcommand{\cf}{cf.\ }
\newcommand{\resp}{respectively\ }
\newcommand{\Subsection}{\subsection}
\providecommand{\nobreakdash}{\nobreak\hbox{-}}
\setlength{\emergencystretch}{2em}
\multlinegap0pt
\usepackage{amsthm}
\usepackage{mathtools}
\usepackage{amsmath,amssymb}
\usepackage{enumitem}
\usepackage{graphicx}
\usepackage{bm}
\usepackage{colonequals}
\newtheorem{thm}{Theorem}[section]
\newtheorem{lem}[thm]{Lemma}
\newtheorem{prop}[thm]{Proposition}
\newtheorem{cor}[thm]{Corollary}
\newtheorem{conj}[thm]{Conjecture}
\theoremstyle{definition}
\newtheorem{defn}[thm]{Definition}
\newtheorem{rmk}[thm]{Remark}
\newtheorem{ex}[thm]{Example}
\newtheorem{quest}[thm]{Question}
\title[Defective chromatic polynomials]{Defective Chromatic Polynomials: the closed form for the linear-forest counts of two cospectral tree families (Lemma 5.2)}
\author[S. Asgarli, T. W. McGinley, N. Xue]{Shamil Asgarli, Tamsen Whitehead McGinley, Nicholas Xue}
\date{Source: arXiv:2605.05550v2 (CC BY 4.0)}
\begin{document}
\maketitle
\vspace{1em}
\noindent\textit{Source and licence.} Defective Chromatic Polynomials. Version arXiv:2605.05550v2 (CC BY 4.0), distributed under the Creative Commons Attribution 4.0 International license, http://creativecommons.org/licenses/by/4.0/, as recorded in the arXiv OAI licence metadata for this exact version and shown on the abstract page https://arxiv.org/abs/2605.05550v2. This item is an editorial re-presentation of the paper's linear-forest expansion lemma together with the paper's own proof of it and the counting lemmas that proof uses. A full rights notice is delivered at licenses/arxiv-2605.05550-CC-BY-4.0.txt.
\noindent\textit{Editorial scope.} Counted unit: the paper's closed form for the linear-forest counts of two cospectral tree families, printed as Lemma 5.2 (source0.tex lines 602--614), delivered together with the paper's complete original proof of it (source0.tex lines 616--692) as the single primary proof body. No mathematics is written, repaired or derived: statement and proof are the paper's own text line for line, in the paper's own notation. Every cross-reference inside the retained spans resolves inside the delivered document, and no citation command occurs in any retained span, so the delivered document carries no bibliography and no reference or citation command of the delivered text was rewritten or adapted. Printed numbers are the paper's own: the wrapper restores the paper's theorem counter before the counted statement, so that the delivered statement prints with the paper's own number. Omitted from this package and not redistributed: the paper's preamble and front matter as such (of which only the title and the author names are reproduced in the wrapper); the introduction, the definitions and the earlier sections; the remaining lemmas, theorems and corollaries; and the paper's bibliography with its bib data file. Printed slips kept literal, among them the paper's own wording and its own choice of symbols. Layout and class: the delivered wrapper is the lane's pinned amsart editorial wrapper at 10pt on A4, to which this package adds the paper's own theorem-like declarations and the shorthand macros its retained text uses, all copied verbatim from the paper's source. No publisher class file, no publisher style file and no upstream script is redistributed. Assets: no retained span of this package contains a figure environment or a graphics-inclusion command, no image or artwork file exists in or is redistributed with this package, and the manifest and the delivery evidence both carry an empty asset-dependency list. A full rights notice is delivered at licenses/arxiv-2605.05550-CC-BY-4.0.txt.
\setcounter{section}{2}
\setcounter{thm}{0}
\begin{thm}\label{thm:general-contraction}
Let $G$ be a graph and let $d\geq 0$. Then
\[
\chi_d(G;k)=\sum_{\substack{A\subseteq E(G)\\ \Delta(G[A])\leq d}} \chi_0(G/A;k).
\]
\end{thm}

\begin{proof}
Given a $(k,d)$-coloring $c$ of $G$, let
\[
A_c\colonequals \{uv\in E(G): c(u)=c(v)\}
\]
be the set of monochromatic edges. Since the number of monochromatic neighbors of a vertex $v$ equals its degree in $G[A_c]$, the defect condition on $c$ is equivalent to $\Delta(G[A_c])\leq d$. Contracting $A_c$, we see that $c$ descends to a coloring of $G/A_c$, since every edge of $A_c$ is monochromatic. No loop arises in $G/A_c$: if an edge $uv\notin A_c$ had both endpoints in the same connected component of $G[A_c]$, then $u$ and $v$ would share a color, forcing $uv\in A_c$, a contradiction. Hence the induced coloring on $G/A_c$ is proper.

Conversely, given $A\subseteq E(G)$ with $\Delta(G[A])\leq d$, each proper coloring of $G/A$ lifts uniquely to a $(k,d)$-coloring of $G$ with $A_c=A$. The lift assigns each vertex $v\in V(G)$ the color of the vertex of $G/A$ containing~$v$. Every edge of~$A$ is then monochromatic by construction. If $uv\notin A$ and $u,v$ were identified by the contraction, $G/A$ would contain a loop, contradicting the existence of a proper coloring; otherwise $u$ and $v$ correspond to distinct adjacent vertices of $G/A$ and receive different colors. Hence no edge outside $A$ is monochromatic.

Thus, the $(k, d)$-colorings of $G$ are partitioned by their monochromatic edge set $A=A_{c}$, and for each admissible $A$, the contribution is $\chi_0(G/A;k)$.
\end{proof}

\setcounter{thm}{1}
\begin{cor}\label{cor:tree-expansion}
Let $T$ be a tree on $n$ vertices and let $d\geq 0$. For each $r\in \{0,1,\dots,n-1\}$, let
\[
c_{r,\leq d}(T)\colonequals \#\{A\subseteq E(T) : |A|=r,\ \Delta(T[A])\leq d\}.
\]
Then
\[
\chi_d(T;k)=\sum_{r=0}^{n-1} c_{r,\leq d}(T)\,k(k-1)^{n-r-1}.
\]
\end{cor}

\begin{proof}
Every subset $A\subseteq E(T)$ gives a forest $T[A]$. Contracting $r=|A|$ of these edges produces a tree on $n-r$ vertices with chromatic polynomial $k(k-1)^{n-r-1}$. Grouping the terms in Theorem~\ref{thm:general-contraction} by $|A|$ gives the result.
\end{proof}

\setcounter{section}{5}
\setcounter{thm}{1}
\begin{lem}\label{lem:linear-forest-counts}
For every $a\geq 1$ and every $r\geq 0$,
\[
c_{r,\leq 2}(X_a)=c_{r,\leq 2}(Y_a)
=
\binom{a+10}{r}
-3\binom{a+7}{r-3}
+\binom{a+5}{r-5}
+2\binom{a+4}{r-6}
-\binom{a+2}{r-8}.
\]
In particular, $X_a$ and $Y_a$ have the same $2$-defective chromatic polynomial.
\end{lem}

\begin{proof}
Set $m=a+10$, the common number of edges. We first count $c_{r,\leq 2}(X_a)$. Since $X_a$ has exactly three vertices of degree $3$, namely $u$, $z_1$, and $w$, an $r$-edge subset of $E(X_a)$ is a linear forest if and only if it avoids selecting all three edges incident with any one of these vertices. For each of these vertices, let $B_u$, $B_{z_1}$, and $B_w$ denote the event that all three incident edges are selected (the ``bad'' events in the inclusion--exclusion). Then
\[
c_{r,\leq 2}(X_a)=\binom{m}{r}-|B_u|-|B_{z_1}|-|B_w|
+|B_u\cap B_{z_1}|+|B_u\cap B_w|+|B_{z_1}\cap B_w|
-|B_u\cap B_{z_1}\cap B_w|.
\]
Each single bad event forces exactly three edges, so
\[
|B_u|=|B_{z_1}|=|B_w|=\binom{m-3}{r-3}.
\]
The edge sets forced by $B_u$ and $B_{z_1}$ share only the edge
$uz_1$, so their union has five edges:
\[
|B_u\cap B_{z_1}|=\binom{m-5}{r-5}.
\]
The events $B_u$ and $B_w$ force disjoint edge sets, so together they force six edges:
\[
|B_u\cap B_w|=\binom{m-6}{r-6}.
\]
Likewise, $B_{z_1}$ and $B_w$ force disjoint edge sets, so
\[
|B_{z_1}\cap B_w|=\binom{m-6}{r-6}.
\]
Finally, the triple intersection forces eight distinct edges:
\[
|B_u\cap B_{z_1}\cap B_w|=\binom{m-8}{r-8}.
\]
Substituting gives
\[
c_{r,\leq 2}(X_a)
=
\binom{m}{r}
-3\binom{m-3}{r-3}
+\binom{m-5}{r-5}
+2\binom{m-6}{r-6}
-\binom{m-8}{r-8}.
\]
Now we count $c_{r,\leq 2}(Y_a)$. Here the three degree-$3$ vertices are $u$, $w$, and $t_1$. As before, let $B_u$, $B_w$, and $B_{t_1}$ denote the corresponding bad events that all three incident edges are selected. Again, an $r$-edge subset is a linear forest if and only if it avoids all three bad events, so
\[
c_{r,\leq 2}(Y_a)=\binom{m}{r}-|B_u|-|B_w|-|B_{t_1}|
+|B_u\cap B_w|+|B_u\cap B_{t_1}|+|B_w\cap B_{t_1}|
-|B_u\cap B_w\cap B_{t_1}|.
\]
Each single bad event again forces three edges:
\[
|B_u|=|B_w|=|B_{t_1}|=\binom{m-3}{r-3}.
\]
The events $B_u$ and $B_w$ force disjoint edge sets, so
\[
|B_u\cap B_w|=\binom{m-6}{r-6}.
\]
The events $B_u$ and $B_{t_1}$ also force disjoint edge sets, so
\[
|B_u\cap B_{t_1}|=\binom{m-6}{r-6}.
\]
The edge sets forced by $B_w$ and $B_{t_1}$ share only the edge
$wt_1$, so their union has five edges:
\[
|B_w\cap B_{t_1}|=\binom{m-5}{r-5}.
\]
Finally, the triple intersection forces eight distinct edges:
\[
|B_u\cap B_w\cap B_{t_1}|=\binom{m-8}{r-8}.
\]
Thus
\[
c_{r,\leq 2}(Y_a)
=
\binom{m}{r}
-3\binom{m-3}{r-3}
+\binom{m-5}{r-5}
+2\binom{m-6}{r-6}
-\binom{m-8}{r-8}.
\]
Hence $c_{r,\leq 2}(X_a)=c_{r,\leq 2}(Y_a)$ for all $a\geq 1$ and $r\geq 0$. The final claim follows from Corollary~\ref{cor:tree-expansion}.
\end{proof}

\end{document}
